\section{Mittelwert, Effektivwert und Leistung}

Diese Section behandelt die \textbf{zentralen Bewertungsgrössen} für periodische, nichtsinusförmige Spannungen und Ströme,
wie sie in Gleichrichtern, Stellern und Wechselrichtern auftreten.

\smallskip

\textbf{Zielbild:}
\begin{itemize}
    \item Definition von Mittelwert und Effektivwert
    \item Physikalische Interpretation (DC-Anteil, thermische Wirkung)
    \item Leistungsberechnung bei allgemeinen periodischen Signalen
    \item Anwendung auf stückweise definierte Zeitverläufe aus Übungen
\end{itemize}

\subsection{Grundlagen: Periodische Signale}

Ein Signal $x(t)$ heisst \textbf{periodisch} mit Periodendauer $T$, wenn gilt:
\[
x(t) = x(t+T) \qquad \forall t.
\]

Typische Beispiele aus der Leistungselektronik:
\begin{itemize}
    \item Gleichgerichtete Sinusspannung (Halb- und Vollwelle)
    \item Strom im R-L-Kreis bei Schaltbetrieb
    \item PWM-Ausgangsspannungen
\end{itemize}

\crd{\textrightarrow\ Mittelwert und Effektivwert werden immer über \textbf{eine volle Periode $T$} gebildet.}

\subsection{Zeitmittelwert (Gleichwert, DC-Anteil)}

\subsubsection{Definition}

Der \textbf{Zeitmittelwert} (Mittelwert) eines periodischen Signals $x(t)$ ist definiert als:
\[
\cbl{\overline{x} = \frac{1}{T} \int_{0}^{T} x(t)\,\diff t.}
\]

\subsubsection{Interpretation und Eigenschaften}

Interpretation:
\begin{itemize}
    \item Entspricht dem \textbf{Gleichanteil} des Signals.
    \item Bei Strom: mittlerer Transport von Ladung pro Zeit.
    \item Bei Spannung: DC-Anteil der Ausgangsspannung eines Gleichrichters.
\end{itemize}

Typische Eigenschaften:
\begin{itemize}
    \item Symmetrische Wechselgrössen: $\overline{x}=0$.
    \item Gleichgerichtete Signale: $\overline{x}>0$.
    \item Für stückweise definierte Signale:
    \[
    \overline{x} = \frac{1}{T} \left( \int_{t_1}^{t_2} x_1(t)\,\diff t + \int_{t_2}^{t_3} x_2(t)\,\diff t + \dots \right).
    \]
\end{itemize}

\crd{\textrightarrow\ Mittelwert ist \textbf{nicht} gleich Effektivwert. Häufige Prüfungsfalle.}

\subsection{Effektivwert (RMS)}

\subsubsection{Definition}

Der \textbf{Effektivwert} eines periodischen Signals $x(t)$ ist definiert als:
\[
\cbl{x_{\mathrm{RMS}} = \sqrt{\frac{1}{T} \int_{0}^{T} x^2(t)\,\diff t}.}
\]

\subsubsection{Interpretation und Vergleich}

Interpretation:
\begin{itemize}
    \item Entspricht dem \textbf{Gleichwert}, der in einem Widerstand \textbf{die gleiche Verlustleistung} erzeugt.
    \item Thermische Wirkung.
    \item Grundlage für Dimensionierung von Leitern, Bauteilen und Verlusten.
\end{itemize}

Vergleich Mittelwert vs. Effektivwert:
\begin{itemize}
    \item Mittelwert: lineare Mittelung.
    \item Effektivwert: quadratische Mittelung.
    \item Allgemein gilt:
    \[
    x_{\mathrm{RMS}} \ge |\overline{x}|.
    \]
\end{itemize}

\subsection{Leistung: Momentanleistung und Wirkleistung}

\subsubsection{Momentanleistung}

Für allgemeine Spannungen und Ströme:
\[
\cbl{p(t) = u(t)\, i(t).}
\]
Diese Grösse schwankt zeitlich, auch im stationären Betrieb.

\subsubsection{Wirkleistung}

Die \textbf{mittlere Wirkleistung} ist definiert als:
\[
\cbl{P = \overline{p(t)} = \frac{1}{T} \int_{0}^{T} u(t)\,i(t)\,\diff t.}
\]

\subsubsection{Spezialfall: Ohmsche Last}

Für $u(t) = R\, i(t)$ gilt:
\[
\cbl{P = R \, I_{\mathrm{RMS}}^2 = \frac{U_{\mathrm{RMS}}^2}{R}.}
\]

Interpretation:
\begin{itemize}
    \item Nur der Effektivwert bestimmt die Verlustleistung.
    \item Unabhängig von der Signalform.
\end{itemize}

\subsection{Rechenworkflow für stückweise definierte Signale}

In vielen Übungen ist $x(t)$ stückweise gegeben:
\[
x(t) =
\begin{cases}
x_1(t), & 0 < t < T_1,\\
x_2(t), & T_1 < t < T_2,\\
\dots
\end{cases}
\]

Dann gilt:

\paragraph{Mittelwert}
\[
\cbl{\overline{x} = \frac{1}{T} \sum_k \int_{T_k} x_k(t)\,\diff t.}
\]

\paragraph{Effektivwert}
\[
\cbl{x_{\mathrm{RMS}} = \sqrt{ \frac{1}{T} \sum_k \int_{T_k} x_k^2(t)\,\diff t }.}
\]

\paragraph{Wirkleistung}
\[
\cbl{P = \frac{1}{T} \sum_k \int_{T_k} u_k(t)\, i_k(t)\,\diff t.}
\]

\crd{\textrightarrow\ Bei stückweisen Verläufen immer korrekt die Zeitintervalle trennen.}

\subsection{Zeitbereichsmodellierung bei Schaltbetrieb (R--L, R--C)}

In der Leistungselektronik entstehen durch Schalter \textbf{stückweise definierte Zeitverläufe}.
In jedem Schaltintervall gilt eine lineare Differentialgleichung niedriger Ordnung.

Ziel:
\begin{itemize}
    \item Zeitverlauf von Strom und Spannung bestimmen,
    \item Einschwingvorgänge analysieren,
    \item Mittelwert, Effektivwert und Ripple berechnen,
    \item stationär-periodischen Betrieb beschreiben.
\end{itemize}

% -------------------------------------------------

\subsubsection{Grundfall: R--L-Last}

Grundgleichung:
\[
\boxed{
u(t) = R\, i(t) + L\, \frac{di(t)}{dt}
}
\]

Dabei:
\begin{itemize}
    \item $u(t)$: angelegte Spannung,
    \item $i(t)$: Laststrom,
    \item $R$: Widerstand,
    \item $L$: Induktivität.
\end{itemize}

Zeitkonstante:
\[
\boxed{
\tau = \frac{L}{R}
}
\]

Physikalische Bedeutung:
\begin{itemize}
    \item Nach $t=\tau$: ca.\ 63\% des Endwerts erreicht.
    \item Nach $t \approx 5\tau$: praktisch eingeschwungen.
\end{itemize}

% -------------------------------------------------

\subsubsection{Lösung bei konstantem Spannungseingang}

Für konstantes $u(t)=U_0$:

\[
\boxed{
i(t) = \frac{U_0}{R} + \left( i(0) - \frac{U_0}{R} \right) e^{-t/\tau}
}
\]

Dabei:
\begin{itemize}
    \item $i(0)$: Anfangsstrom,
    \item $\frac{U_0}{R}$: stationärer Endstrom.
\end{itemize}

Interpretation:
\begin{itemize}
    \item Exponentielles Einschwingen.
    \item Geschwindigkeit durch $\tau$ bestimmt.
\end{itemize}

% -------------------------------------------------

\subsubsection{Kontinuitätsbedingungen bei Schaltvorgängen}

Bei idealen Bauteilen gilt:

\[
\boxed{i_L(t^-) = i_L(t^+)}
\qquad
\boxed{u_C(t^-) = u_C(t^+)}
\]

Bedeutung:
\begin{itemize}
    \item Induktorstrom ist stetig.
    \item Kondensatorspannung ist stetig.
    \item Diese Bedingungen liefern die Anfangswerte für jedes neue Intervall.
\end{itemize}

% -------------------------------------------------

\subsubsection{Stückweise Betrieb bei Schaltperioden}

Typisch im Chopper oder PWM:

EIN-Intervall:
\[
L \frac{di}{dt} + R i = U_{DC}
\]

AUS-Intervall:
\[
L \frac{di}{dt} + R i = 0
\]

Lösungen:

EIN:
\[
\boxed{
i(t) = \frac{U_{DC}}{R} + \left( i_{\min} - \frac{U_{DC}}{R} \right) e^{-(t-t_e)/\tau}
}
\]

AUS:
\[
\boxed{
i(t) = i_{\max} \, e^{-(t-t_a)/\tau}
}
\]

Dabei:
\begin{itemize}
    \item $t_e$: Einschaltzeitpunkt,
    \item $t_a$: Ausschaltzeitpunkt,
    \item $i_{\min}$, $i_{\max}$: Stromgrenzen pro Periode.
\end{itemize}

% -------------------------------------------------

\subsubsection{Stationär-periodischer Betrieb}

Im eingeschwungenen Schaltbetrieb gilt:

\[
\boxed{
i(t) = i(t + T_s)
}
\]

mit:
\begin{itemize}
    \item $T_s$: Schaltperiode.
\end{itemize}

Diese Periodenbedingung erlaubt die Bestimmung von $i_{\min}$ und $i_{\max}$.

% -------------------------------------------------

\subsubsection{Stromrippel}

Definition:
\[
\boxed{
\Delta i = i_{\max} - i_{\min}
}
\]

Näherung bei grossem $L$:
\[
\boxed{
\Delta i \approx \frac{(U_{DC} - \overline{u}) \, d T_s}{L}
}
\]

Dabei:
\begin{itemize}
    \item $d$: Tastgrad,
    \item $\overline{u}$: gemittelte Ausgangsspannung.
\end{itemize}

% -------------------------------------------------

\subsubsection{Mittelwert des Stroms}

\[
\boxed{
I_{\mathrm{av}} = \frac{1}{T_s} \int_0^{T_s} i(t)\, dt
}
\]

Für kleines Ripple:
\[
\boxed{
I_{\mathrm{av}} \approx \frac{i_{\max} + i_{\min}}{2}
}
\]

% -------------------------------------------------

\subsubsection{Physikalische Kernaussagen}

\begin{itemize}
    \item Zeitkonstante bestimmt Einschwinggeschwindigkeit.
    \item Kontinuität bestimmt Anfangswerte nach jedem Schalten.
    \item Stationär bedeutet periodisch, nicht konstant.
    \item Mittelwert bestimmt Drehmoment, Leistung, Arbeitspunkt.
    \item Ripple bestimmt Verluste und Dimensionierung.
\end{itemize}

\subsection{Bezug zu Fourier-Reihen}

Für ein periodisches Signal mit Fourierkoeffizienten gilt (Parseval):
\[
\cbl{X_{\mathrm{RMS}}^2 = \frac{a_0^2}{2} + \sum_{n=1}^\infty \left( a_n^2 + b_n^2 \right).}
\]

Interpretation:
\begin{itemize}
    \item Effektivwert ist die Summe der Quadrate aller Harmonischen.
    \item Harmonische tragen direkt zur Verlustleistung bei.
\end{itemize}


\subsection{Skizze : Leistung bei ohmscher Last}

\begin{center}
\begin{tikzpicture}[scale=1]
\draw[->] (0,0) -- (8,0) node[right] {$t$};
\draw[->] (0,0) -- (0,3) node[above] {$p(t)$};

\draw (0,1.5) .. controls (1,2.5) .. (2,1.5)
      .. controls (3,0.5) .. (4,1.5)
      .. controls (5,2.5) .. (6,1.5)
      .. controls (7,0.5) .. (8,1.5);

\draw[dashed] (0,1.5) -- (8,1.5) node[right] {$P$};
\end{tikzpicture}
\end{center}

\textbf{Aussage:}
\begin{itemize}
    \item Momentanleistung schwankt.
    \item Wirkleistung ist der Mittelwert.
\end{itemize}

\subsection{Typische Fehlerquellen}

\begin{itemize}
    \item Mittelwert mit Effektivwert verwechselt.
    \item Quadrat beim RMS vergessen.
    \item Falsche Periodendauer $T$ eingesetzt.
    \item Stückweise Integrale nicht getrennt.
    \item Leistung mit $U_{\mathrm{RMS}}\cdot I_{\mathrm{RMS}}$ berechnet bei nichtohmscher Last.
\end{itemize}

\subsection{Minimal-Checkliste für Aufgaben}

\begin{enumerate}
    \item Periodendauer $T$ bestimmen.
    \item Signalform stückweise festlegen.
    \item Mittelwert berechnen.
    \item Effektivwert berechnen.
    \item Bei Leistung: $p(t)=u(t)i(t)$ bilden und mitteln.
    \item Plausibilität prüfen: $x_{\mathrm{RMS}} \ge |\overline{x}|$.
\end{enumerate}

